\begin{center}\bfseries Abstract\end{center}
\begin{narrowtext}
Young sub-Neptunes may hold magma that is fully miscible with their hydrogen envelopes. On
cooling, hydrogen exsolves and releases heat that none of the evolution models we checked
includes. We add energy, nucleation and cooling-path functions to the
open-source code \Atm{} and apply them to the 2026 H$_2$--MgSiO$_3$ mixing model and its
source simulations.\par
\vspace{3pt}
At 1--3\,wt\% hydrogen and 2--5\,GPa, exsolution releases \hbarMOneLo{}--\hbarMOneHi{}\,MJ per
kg of hydrogen in refits to the simulated excess enthalpies, and
\hbarPubLo{}--\hbarPubHi{}\,MJ\,kg$^{-1}$ in the published fit, which underestimates the
simulated excess enthalpy of the most hydrogen-poor mixture (\edHexPub{} against
$\edHexData\pm\edHexSig$\,kJ\,mol$^{-1}$). A convecting interior degasses into the
envelope without nucleation. At 3\,wt\% and 2\,GPa, its exsolution heat over the first 300\,K
below the binodal is \intRatioThree{} (\intRatioThreeLo{}--\intRatioThreeHi{}) times the
sensible heat, and its effective heat capacity rises up to \intCeffThree{}-fold if the hydrogen
leaves as pure H$_2$. Scaled onto a published evolution built on an earlier phase diagram, the
heat is \planetRadiatedLo{}--\planetRadiatedHi{}\% of the energy radiated.\par
\vspace{3pt}
Away from the envelope, homogeneous nucleation with the simulated interface width ends
metastability \dlyFigOneCPubLo{}--\dlyFigOneCPubHi{}\,K below the binodal (refit:
\dlyFigOneCMOneLo{}--\dlyFigOneCMOneHi{}\,K) and withholds at most \whFigOneCPubHi{}\% of the
hydrogen. The \pubIntTwo{}\,K binodal--spinodal window of the published fit is thus an
unreached mean-field bound. Refits give median windows of \IntRangeLo{}--\IntRangeHi{}\,K. A
10\% curvature measurement at \curvT{}\,K would lower the standard deviation of the spinodal
only from \designSdNow{} to \curvSdPooled{}\,K, and only through the assumed model forms.
\end{narrowtext}
\vspace{10pt}
\noindent\rule{\textwidth}{0.6pt}
\vspace{10pt}
\begin{center}\bfseries Samenvatting\end{center}
\begin{narrowtext}
Jonge sub-Neptunussen kunnen magma bevatten dat volledig mengbaar is met hun
waterstofomhulsel. Bij afkoeling ontmengt waterstof, en de vrijgekomen warmte ontbreekt in
alle evolutiemodellen die wij nagingen. Wij voegen functies voor energie,
kiemvorming en afkoelingspaden toe aan de opensourcecode \Atm{} en passen ze toe op het
mengmodel voor H$_2$--MgSiO$_3$ uit 2026 en op de simulaties waarop het berust.\par
\vspace{3pt}
Bij 1--3\,massa\% waterstof en 2--5\,GPa maakt ontmenging \hbarMOneLo{}--\hbarMOneHi{}\,MJ per
kg waterstof vrij in nieuwe fits aan de gesimuleerde exces-enthalpie\"en, en
\hbarPubLo{}--\hbarPubHi{}\,MJ\,kg$^{-1}$ in de gepubliceerde fit, die de gesimuleerde
exces-enthalpie van het waterstofarmste mengsel onderschat (\edHexPub{} tegenover
$\edHexData\pm\edHexSig$\,kJ\,mol$^{-1}$). Een convectief binnenste ontgast zonder
kiemvorming in het omhulsel. Bij 3\,massa\% en 2\,GPa is zijn ontmengingswarmte over de eerste
300\,K onder de binodaal \intRatioThree{} (\intRatioThreeLo{}--\intRatioThreeHi{}) keer de
voelbare warmte, en stijgt zijn effectieve warmtecapaciteit tot het \intCeffThree{}-voudige als
de waterstof als zuiver H$_2$ vertrekt. Geschaald op een gepubliceerde evolutie met een ouder
fasediagram is de warmte \planetRadiatedLo{}--\planetRadiatedHi{}\% van de uitgestraalde
energie.\par
\vspace{3pt}
Ver van het omhulsel be\"eindigt homogene kiemvorming met de gesimuleerde grensvlakbreedte de
metastabiliteit \dlyFigOneCPubLo{}--\dlyFigOneCPubHi{}\,K onder de binodaal (nieuwe fit:
\dlyFigOneCMOneLo{}--\dlyFigOneCMOneHi{}\,K), en blijft hoogstens \whFigOneCPubHi{}\% van de
waterstof achter. Het binodaal-spinodaalvenster van \pubIntTwo{}\,K in de gepubliceerde fit
is dus een onbereikte gemiddeld-veldbovengrens. Nieuwe fits geven mediane
vensters van \IntRangeLo{}--\IntRangeHi{}\,K. Een krommingsmeting met 10\% precisie bij
\curvT{}\,K zou de standaardafwijking van de spinodaal slechts van \designSdNow{} naar
\curvSdPooled{}\,K verkleinen, en alleen via de aangenomen modelvormen.
\end{narrowtext}
